% !TeX root = ../papers/nc-p3.tex
\section{反対根の二次接触と少なくとも十台点での八次判定}
\label{sec:opposite-contact-support}
\begin{theorem}\label{thm:opposite-contact-support}
四辺有限な標準有限Frobenius窓で、配置 $X$ の台が十点以上なら、
射影非同値な反対根は必ず $H^0(\mathcal I_X^3(8))$ に非零の八次式を与える。
特にこの空間が零の有限辺からの標準正則延長は全成分を通じて高々一つで、
転置辺を使う有限検査で存在を判定できる。
\end{theorem}

\begin{proof}[接触の消去イデアルと平面的な欠損]
行列式が恒等零と仮定する。前節の定理でその余因子を $ff^{\mathsf t}$ と書き、
$f_i$ は四つの独立な三次式である。三次評価は同型なので $X$ の関数環における
これらの線型包も四次元である。各点 $p$ の局所環を $D$、極大イデアルを $\mathfrak m$、
長さを $n$ とし、$f_i$ の生成する局所イデアルを $J$ とする。
前節と同じ単元成分でSchur補を取る。一般階数2で二次小行列式が単元平方を除いて平方なので、
完備正則局所環上で
\[
           S=q(A B^{\mathsf t}+B A^{\mathsf t}),\qquad
           f'=\text{単元}\cdot q(A\times B)
\]
と書ける。$f'$ は最後の三成分で、第一成分はその局所的な線型結合である。
今は共通因子 $q$ を許す。$K$ を $A,B$ の全成分のイデアルとする。
$S|_D=0$ で $q A_iA_jB_k=q A_iB_jB_k=0$ となり $KJ=0$。
また各 $f_i$ の全一階偏導は $K$ に入る。これらの主張は任意の局所自明化で成り立つ。
$J\ne0$ なら $A,B$ はともに消え、$S$ は二次以上である。
前節の三次接触の排除により $q$ は単元で、$A,B$ はそれぞれ非零一次部を持つ。

一方、$f_i^2\in I_p^2$ より $f_i\in\overline I_p$ である。
埋込次元2以下の局所点での評価の次元は
\[
       \delta(n)=\left\lfloor (n-\sqrt n)/2\right\rfloor
\]
以下である。曲線状では既出の整閉性により零とする。平面的な場合は滑らかな一変数を
消去して二生成パラメータ・イデアルの整閉包を取り、その重複度は $n$ である。
二変数有限余長イデアルの単項式初期イデアル $M$ は余長を保ち、重複度 $e(M)$ は元以上である。
Newton多角形の両軸切片を $a,b$、コンパクト辺の格子長の和を $l$ とする。
余長は折れ線より原点側の格子点数以上で、Pickの公式から
$2\operatorname{length}(R/M)\ge e(M)+a+b-l$ である（ここだけ $R$ は二変数環）。
各格子長は辺の各座標差以下なので $l\le\min(a,b)$。
さらに折れ線は両切片を結ぶ線分以下なので $e(M)\le ab$ であり、従って
$2\operatorname{length}(R/M)\ge n+\sqrt n$。これで増分の上界が従う。
全点が平面的なら、台が十点以上のもとで $\sum(n-1)\le10$ より $\sum\delta(n)\le3$。
実際 $n=1,2,3$ で零、$4,5,6$ で1、$7,8$ で2、$9,10,11$ で3の表から従う。
四次元に反する。
\end{proof}

\begin{proof}[非平面点での二次部]
埋込次元3の点は長さ8以上なので今は高々一点 $p$ で、その長さは11以下である。
三つの局所完全交叉生成元の二次部は独立である。従属なら一生成元の位数を3以上にでき、
局所交点重複度の積下界から $n\ge2\cdot2\cdot3=12$ となる。
この下界は、対応する位数の一般斉次式を各生成元にパラメータ倍して加えれば
一般の初期形式が正則列となり、長さの上半連続性で特殊点の長さがその次数積以上となることでも従う。
二次部の三空間を $Q$ とする。特に $\operatorname{length}\mathfrak m^3=n-7$。

$p$ における四本の三次式の二次部の線型包を $E$ とする。定数部と一次部は零であり、
$E^2\subseteq Q^2$ である。ここで積は三変数多項式環の内部で取る。
$\dim E\ge2$ なら $E\subset Q$ となることを示す。
まず $Q$ の三生成元が代数的に独立なら、その多項式環の中で各 $e_i e_j$ を
$Q$ の生成元について二次式として書く。この対称行列は階数1であり、一意分解により
$h v_i v_j$ と斉次に書ける。$E$ の独立性から $v_i$ は全て定数にはできず、
従って $h$ は定数で各 $v_i$ が一次式である。対角項を元の環へ戻せば $e_i\in Q$。
$Q$ に共通一次因子があれば $Q$ はその一次因子と全一次式との積であり、
$e_i^2\in Q^2$ からやはり従う。共通因子がなく代数的従属ならその射影像は非退化曲線で、
一般直線からの像の次数で高々2、よって滑らかな二次曲線である。座標変更で
$q_1q_3=q_2^2$ とでき、原始性から $Q$ は二つの一次式の全二次式になる。
この二変数だけの多項式が平方ならその平方根もそこに属し、主張が従う。

従って $E$ の次元にかかわらず、局所での四本の線型包の次元は高々 $1+(n-7)$ である。
$n=8$ なら他点で残る余分な長さは3以下なので、それらの寄与は高々1。
$n=9$ なら他点の寄与は零である。いずれも四次元には達しない。
残る $n=10,11$ では他点は全て曲線状であり、四本は $p$ の局所環だけで独立である。
\end{proof}

\begin{proof}[中心射影と局所socle]
まず $\dim E\le1$ とする。少なくとも三本の独立な三次式を線型結合で $p$ に三重に消せる。
これらは $p$ を頂点とする三次錐であり、方向平面上の三次式である。
$p$ 以外の少なくとも九台点の中心射影は相異なる。$p$ を通る直線の $p$ での
切断長は2以上であり、直線全体の切断長は3以下だからである。
また射影平面で一直線上の像は高々四点である。その逆像平面での $p$ の切断長は
埋込次元から3以上、全切断長は7以下だからである。

平面上の三本の独立な三次式がこのような九点以上に消えれば共通二次曲線を持つ。
固定成分がなければ一般二本の交点は長さ9であり、九点を含むならその被約完全交叉に等しい。
その三次イデアルは二次元なので第三の独立な式はない。固定直線を除いた場合も、
残りの少なくとも五点で三本の二次式が消えるため、次数積4からさらに共通因子がある。
従って三本は一つの二次式 $C$ と三つの独立な一次式との積である。
全ての像点は $C$ に乗り、二直線では八点までなので $C$ は非特異である。
$p$ でのアフィン座標を $w_1,w_2,w_3$ とすれば、三本は $C(w)w_i$ に取れる。
その一階偏導の包は $w$ の全二次式である。非特異二次式を平方和に取ると
非対角積が出て、対角部の係数行列も $2I+\boldsymbol{1}\boldsymbol{1}^{\mathsf t}$ で可逆である。
これらはすべて $J$ を消すので $J\subseteq\operatorname{ann}\mathfrak m^2$。
右辺はGorenstein双対性により長さ4であり、独立な四評価から等しい。
$KJ=0$ より $K D\subseteq\mathfrak m^2$ である。
埋込次元3なのでこれは $A,B$ の全成分の一次部が零であることを意味し、二次接触に反する。

最後に $\dim E\ge2$ なら上で示した $E\subset Q$ から四本は $D$ で全て $\mathfrak m^3$ に入る。
$n=10$ ではこれは長さ3しかなく、$n=11$ では独立性から $J=\mathfrak m^3$ となる。
$A,B$ の一次部の全成分が張る空間の次元は2以上である。そうでなければ
その外積の二次部がただ一つの平方の倍に限られ、第一核座標も $p$ で零の係数による
他成分の和なので $\dim E\le1$ となる。よって $D/KD$ の埋込次元は高々1。
$J$ はこの一変数Artin環の加群でsocleが一次元以下である（$D$ のsocleに含まれる）。
一変数の冪零作用のJordan分解から $J$ は巡回で、$\dim\mathfrak m^3/\mathfrak m^4=1$。

一方、$A,B$ の一次部の成分の任意の積（一方ずつ取る）は $Q$ に入る。
対称積は関係 $S=0$ の二次部、外積は $E\subset Q$ で、$q$ は単元だからである。
少なくとも一方の列の成分は二つの独立な一次式を含む。両方がそれぞれ一式の倍だけなら
外積は一つの二次式の倍になる。他方の非零一次式をこれら二つに掛けると、
$Q$ の二つの独立な元の間に線型シジジーを得る。
局所三生成元の二次部は独立だから、イデアルの三次初期部分は $S_1Q$ そのものであり、
このシジジーにより次元は8以下である。三変数の三次式は十次元なので
$\dim\mathfrak m^3/\mathfrak m^4\ge2$ となり矛盾する。
以上で行列式は非零である。三重イデアルへの所属は単元成分での消去から、
延長の主張は系\ref{cor:octic-free-finite-extension}と四辺の有限性から従う。
\end{proof}
